%! AP Statistics — Unit 8, Lesson 8.4 — Lesson Plan
\documentclass[10pt]{article}
\usepackage{apstats-article}
\usepackage{apstats-boxes}
\newcommand{\TallMath}[1]{$\displaystyle #1\rule[-1.4em]{0pt}{3.2em}$}
\newcommand{\UnitNumberName}{Unit 8: Inference for Categorical Data: Chi-Square \quad}
\newcommand{\LessonNumberName}{Lesson 8.4: Expected Counts in Two-Way Tables}

\begin{document}
\begin{center}
    {\Large\bfseries \CourseName: \SchoolYear \\
    {\small\bfseries \UnitNumberName \LessonNumberName}}
\end{center}

\begin{tcolorbox}[colback=sky, colframe=navy, boxrule=0.9pt, arc=2mm,
    left=2mm, right=2mm, top=1.5mm, bottom=1.5mm]
    \textbf{Primary Objective:} Students will calculate expected counts in two-way tables
    and explain what those counts represent under $H_0$ --- specifically, what cell counts
    would look like if the two categorical variables were independent (or the populations
    were homogeneous). (Big Idea: VAR)
\end{tcolorbox}

\renewcommand{\arraystretch}{1.6}
\begin{skillbox}[Priority Ideas \& Skills]{goldbox}
\begin{minipage}[t]{0.48\linewidth}
    \textbf{AP Skill 3.A --- Determining Parameters, Statistics, and Sampling Distributions}
    \begin{itemize}
        \item Calculate expected counts in a two-way table using the formula
              $\text{Expected} = \dfrac{(\text{row total})(\text{col total})}{\text{table total}}$
              (LO VAR-8.H; Skill 3.A).
        \item Derive the formula from two equivalent perspectives:
              (row proportion)(column total) \emph{or} (row total)(column proportion) ---
              same result either way.
        \item Bridge lesson: connects GOF (Lesson 8.3) to chi-square tests for
              homogeneity and independence (Lessons 8.5--8.6).
    \end{itemize}
\end{minipage}
\hfill
\begin{minipage}[t]{0.48\linewidth}
    \textbf{Key Understandings}
    \begin{itemize}
        \item Expected counts represent the cell frequencies we would \emph{expect} if the
              two variables were completely independent --- or if the population distributions
              were identical across groups.
        \item The large-counts condition for two-way chi-square tests: all expected counts
              must be at least 5.
        \item Equal row (or column) totals do \emph{not} guarantee equal expected counts ---
              expected counts also depend on the other marginal totals.
        \item This lesson is a computation bridge; the full chi-square test (comparing
              observed vs.\ expected) is built in Lessons 8.5--8.6.
    \end{itemize}
\end{minipage}
\end{skillbox}

\begin{skillbox}[{Vocabulary, Concepts, \& Theorems}]{greenbox}
\begin{minipage}[t]{0.48\linewidth}
    \begin{itemize}
        \item \textbf{Two-way table} --- a display of counts for two categorical variables;
              rows = one variable, columns = the other
        \item \textbf{Row total} --- sum of counts across one row (one category of the
              row variable)
        \item \textbf{Column total} --- sum of counts down one column (one category of the
              column variable)
    \end{itemize}
\end{minipage}
\hfill
\begin{minipage}[t]{0.48\linewidth}
    \begin{itemize}
        \item \textbf{Expected count (two-way)} --- the count we would predict in a cell
              if the row and column variables were independent:
              $E = \dfrac{(\text{row total})(\text{col total})}{\text{table total}}$
        \item \textbf{Marginal distribution} --- the row totals or column totals, expressed
              as counts or proportions; describes one variable ignoring the other
    \end{itemize}
\end{minipage}
\end{skillbox}

\begin{skillbox}[Activate Prior Knowledge \& Spiral Review]{sky}
\begin{minipage}[t]{0.48\linewidth}
    \textbf{Prior Knowledge Check (3--4 min)}
    \begin{itemize}
        \item Lesson 8.3: chi-square GOF statistic formula and interpretation.
        \item Lesson 8.3: expected counts for a one-way table (GOF) --- expected $= n \cdot p_0$.
        \item Lesson 8.1--8.2: purpose of chi-square tests; conditions (random, large counts,
              independence).
        \item Connect: ``In GOF we compared observed to expected for one variable. Today
              we extend that idea to \emph{two} variables simultaneously.''
    \end{itemize}
\end{minipage}
\hfill
\begin{minipage}[t]{0.48\linewidth}
    \textbf{Warm-Up (attached)}\\
    Spiral review of chi-square GOF:
    \begin{itemize}
        \item Decision and conclusion from a given $\chi^2$ statistic/p-value.
        \item Interpreting degrees of freedom in context.
    \end{itemize}
    \textbf{Connections Forward}
    \begin{itemize}
        \item Chi-square test for homogeneity $\to$ Lesson 8.5
        \item Chi-square test for independence $\to$ Lesson 8.6
    \end{itemize}
\end{minipage}
\end{skillbox}

\begin{skillbox}[Hook]{sky}
\textbf{Motivating Question (2--3 min)}\\
Display the grade/learning-preference table (200 students). Ask:
\textbf{``If grade level and learning preference are completely unrelated, how many
9th-graders would we \emph{expect} to prefer online learning?''}

Let students guess, then reveal the formula connection: if 50 out of 200 are 9th-graders
(25\% of the sample) \emph{and} 100 out of 200 prefer online (50\% of the sample),
independence predicts $0.25 \times 100 = 25$ students in that cell.
That is exactly $(50 \times 100)/200 = 25$.
\end{skillbox}

\begin{skillbox}[Lesson]{sky}
\begin{minipage}[t]{0.48\linewidth}
    \textbf{Part 1 --- Reading a Two-Way Table (5 min)}
    \begin{enumerate}
        \item Display the grade/learning table. Label: rows = grade level, columns = preference,
              cells = joint counts, margins = row and column totals.
        \item Ask students to compute each marginal proportion.
              Emphasize: the marginal distribution of preference is 100/200 = 50\% online,
              50\% in-person --- regardless of grade.
        \item Pose: ``Under independence, what fraction of any grade level should prefer
              online?'' Answer: 50\% $\to$ so expected count = $0.50 \times 50 = 25$.
    \end{enumerate}

    \textbf{Part 2 --- Deriving the Formula (5 min)}\\
    Two equivalent routes to the same answer:
    \begin{itemize}
        \item \textit{Route A}: (row proportion)(column total)
              $= \dfrac{50}{200} \times 100 = 25$
        \item \textit{Route B}: (row total)(column proportion)
              $= 50 \times \dfrac{100}{200} = 25$
        \item Both simplify to $\dfrac{(\text{row total})(\text{col total})}{\text{table total}}$.
    \end{itemize}
\end{minipage}
\hfill
\begin{minipage}[t]{0.48\linewidth}
    \textbf{Part 3 --- Filling in the Full Table (5 min)}

    \renewcommand{\arraystretch}{1.4}
    \begin{tabular}{lccc}
        & \textbf{Online} & \textbf{In-person} & \textbf{Total} \\
        \hline
        9th  & 28 & 22 & 50 \\
        10th & 35 & 15 & 50 \\
        11th & 22 & 28 & 50 \\
        12th & 15 & 35 & 50 \\
        \hline
        Total & 100 & 100 & 200 \\
    \end{tabular}

    \vspace{0.1in}
    All expected counts $= (50 \times 100)/200 = 25$
    (balanced design; each row is 25\% of total; each column is 50\%).

    \vspace{0.1in}
    \textbf{Large-Counts Condition:}\\
    All expected counts $\ge 5$? \quad $25 \ge 5$ \checkmark\\
    (Condition met for all 8 cells.)

    \vspace{0.1in}
    \textbf{Key Observation:} The observed counts differ from 25 in some cells ---
    this discrepancy is what the chi-square statistic will measure in Lesson 8.5.
\end{minipage}
\end{skillbox}

\begin{skillbox}[Active Monitoring]{redbox}
\begin{minipage}[t]{0.48\linewidth}
    \textbf{Circulate and Check}
    \begin{itemize}
        \item Students often apply the formula to cell counts instead of row/column totals.
              Redirect: ``Which number comes from the row margin? Which from the column margin?''
        \item Watch for students dividing row total by table total, then stopping --- remind
              them to multiply by the column total.
        \item Check that students distinguish \emph{observed} counts (given in the table)
              from \emph{expected} counts (calculated).
    \end{itemize}
\end{minipage}
\hfill
\begin{minipage}[t]{0.48\linewidth}
    \textbf{Cold-Call Prompts}
    \begin{itemize}
        \item ``Without computing, what would all expected counts look like if every
              row total were equal \emph{and} every column total were equal?'' (They'd
              all be the same.)
        \item ``If row totals are unequal but column totals are equal, are expected
              counts in a given column all the same?'' (No --- they scale with row totals.)
        \item ``What does it mean for the large-counts condition to fail?''
    \end{itemize}
\end{minipage}
\end{skillbox}

\begin{skillbox}[Group Work \& Differentiation]{redbox}
\begin{minipage}[t]{0.30\linewidth}
    \textbf{Tier R --- Remediate}
    \begin{itemize}
        \item Provide the formula card: $E = RT \times CT / N$.
        \item Supply a partially completed expected-count table; students fill the
              remaining cells.
        \item Sentence frame for condition check: ``All expected counts are
              \underline{\hspace{1cm}}, which is [greater than / less than] 5, so
              the large-counts condition is [met / not met].''
    \end{itemize}
\end{minipage}
\hfill
\begin{minipage}[t]{0.30\linewidth}
    \textbf{Tier A --- Approaching}
    \begin{itemize}
        \item Complete the full activity: compute all 6 expected counts in the sports
              participation table, check the large-counts condition, and compare observed
              vs.\ expected.
        \item Show work for at least 2 cells using the formula explicitly.
    \end{itemize}
\end{minipage}
\hfill
\begin{minipage}[t]{0.30\linewidth}
    \textbf{Tier E --- Extension}
    \begin{itemize}
        \item After computing expected counts, identify the grade/sport combination
              showing the largest deviation from expected.
        \item Conjecture: does this suggest an association between grade and sports
              participation? Explain using both numbers and words.
        \item Preview: write the $\chi^2$ contribution for that one cell:
              $(O - E)^2 / E$.
    \end{itemize}
\end{minipage}
\end{skillbox}

\begin{skillbox}[Individual Work \& Assessment]{redbox}
\begin{minipage}[t]{0.48\linewidth}
    \textbf{Exit Ticket (5 min)}\\
    A 2$\times$3 table (120 students; gender vs.\ favorite subject):
    \begin{itemize}
        \item Male: Math 22, Science 28, English 10 $\to$ row total 60
        \item Female: Math 18, Science 17, English 25 $\to$ row total 60
        \item Column totals: 40, 45, 35; Table total: 120
    \end{itemize}
    Students calculate:
    \begin{enumerate}
        \item Expected count for Male/Math: $(60 \times 40)/120 = 20$.
        \item Expected count for Female/English: $(60 \times 35)/120 = 17.5$.
    \end{enumerate}
\end{minipage}
\hfill
\begin{minipage}[t]{0.48\linewidth}
    \textbf{Collection Note}\\
    Collect exit tickets and sort into three piles:
    \begin{itemize}
        \item Correct formula and arithmetic --- ready for Lesson 8.5.
        \item Correct formula, arithmetic error --- reteach before Lesson 8.5.
        \item Incorrect formula --- pull for small-group review of marginals.
    \end{itemize}
    \textbf{AP-Style Check:} ``True or False: If all expected counts in a two-way
    table are greater than 5, the observed counts must also be greater than 5.''
    (False --- the large-counts condition only requires expected $\ge 5$.)
\end{minipage}
\end{skillbox}

\begin{skillbox}[Reinforcement \& Extension]{goldbox}
\begin{minipage}[t]{0.48\linewidth}
    \textbf{Homework}
    \begin{itemize}
        \item Problem 1: compute all 6 expected counts in a 3$\times$2 table of
              240 students (grade 10/11/12 vs.\ car ownership).
        \item Problem 2: True/False --- ``equal row totals imply equal expected counts.''
              Justify.
        \item Problem 3: 3$\times$4 table, $n = 300$; all row totals $= 100$, all column
              totals $= 75$. Compute expected counts; check large-counts condition.
    \end{itemize}
\end{minipage}
\hfill
\begin{minipage}[t]{0.48\linewidth}
    \textbf{Extension / Preview}
    \begin{itemize}
        \item \textit{Extension}: For the grade/learning table, compute $(O - E)^2/E$
              for each of the 8 cells and sum them. What number do you get? (This is the
              $\chi^2$ statistic students will use in Lesson 8.5.)
        \item \textit{Preview}: Lesson 8.5 introduces the chi-square test for homogeneity ---
              the full hypothesis-testing procedure that uses the expected counts computed
              today.
    \end{itemize}
\end{minipage}
\end{skillbox}

\end{document}
